\documentclass[12pt]{article}
\usepackage[utf8]{inputenc}
\usepackage[T1]{fontenc}
\usepackage[margin=1in]{geometry}
\usepackage{setspace}
\usepackage{xcolor}
\usepackage[colorlinks=true, urlcolor=blue!60!black, linkcolor=blue!60!black]{hyperref}
\setlength{\parindent}{2em}
\setlength{\parskip}{1.25em}
\renewcommand{\baselinestretch}{1.0}
\title{\textit{Shaw v. Reno}, 509 U.S. 630 (1993)}
\author{Case Brief}
\date{}
\begin{document}
\maketitle

\section{Case Metadata}
\begin{description}
  \item[Case Name:] \textit{Shaw v. Reno}, 509 U.S. 630 (1993) (No. 92-357).
  \item[Decided:] June 28, 1993 (Argued April 20, 1993).
  \item[Vote:] 5--4 (reversed and remanded).
  \item[Author:] Justice Sandra Day O'Connor (joined by Chief Justice Rehnquist and Justices Scalia, Kennedy, and Thomas).
  \item[Dissents:] Justice White (joined by Justices Blackmun and Stevens); Justice Blackmun; Justice Stevens; Justice Souter.
  \item[Chapter / Topic:] Chapter 7: Equal Protection --- Racial Gerrymandering and Voting Rights.
  \item[Primary Source URLs:] \href{https://supreme.justia.com/cases/federal/us/509/630/}{Justia (509 U.S. 630)}; \href{https://supreme.justia.com/cases/federal/us/509/630/case.pdf}{Justia Opinion PDF}; \href{https://www.law.cornell.edu/constitution-conan/amendment-15/racial-gerrymandering-and-right-to-vote-clause}{Cornell LII (Constitution Annotated)}.
\end{description}

\section{Facts and Procedural History}
Following the 1990 census, North Carolina---whose voting-age population was roughly 78\% white and 20\% Black, with Black residents dispersed across the State and forming a majority in only 5 of 100 counties---gained a twelfth seat in the U.S. House of Representatives. Because 40 of its 100 counties were covered jurisdictions under \S~5 of the Voting Rights Act of 1965 (VRA), North Carolina submitted its initial congressional redistricting plan, which included one majority-Black district, to the U.S. Attorney General for preclearance. The Attorney General formally objected, asserting that the General Assembly could have drawn a second majority-minority district in the south-central to southeastern region.

Rather than seeking a declaratory judgment in the D.C. District Court, the General Assembly enacted a revised plan creating two majority-Black districts: a hook-shaped District 1 in the northeast (likened to a ``Rorschach ink-blot test'') and an extraordinarily irregular District 12 in the north-central Piedmont. District 12 snaked approximately 160 miles along the Interstate 85 corridor---often no wider than the highway itself---splitting 10 counties and numerous towns to link dispersed enclaves of Black neighborhoods; one legislator quipped that driving down I-85 with both car doors open would hit most of the district's residents. The Attorney General precleared the revised plan.

Five white residents of Durham County, led by Ruth Shaw (two placed in District 12 and three in neighboring District 2), sued state and federal officials in the Eastern District of North Carolina. They did not allege white vote dilution; instead, they claimed that deliberately segregating voters into bizarrely shaped districts by race violated their Fourteenth Amendment right to a ``color-blind'' electoral process. A three-judge District Court unanimously dismissed the federal defendants for lack of jurisdiction and dismissed the Equal Protection claim against the state defendants (2--1), reading \textit{United Jewish Organizations of Williamsburgh, Inc. v. Carey}, 430 U.S. 144 (1977) (\textit{UJO}), to bar white voters' claims absent discriminatory purpose and statewide proportional underrepresentation. Chief Judge Voorhees dissented in part, and the Supreme Court noted probable jurisdiction.

\section{Issue(s) Presented}
Whether voters state a cognizable, analytically distinct claim under the Equal Protection Clause of the Fourteenth Amendment---triggering strict scrutiny---when they allege that a facially race-neutral redistricting plan is so bizarrely shaped that it is unexplainable on grounds other than race and subordinates traditional race-neutral districting principles to racial segregation, even without alleging dilution of any racial group's voting strength.

\section{Holding}
Yes (5--4). Appellants stated an analytically distinct Equal Protection claim by alleging that North Carolina's reapportionment scheme was so irrational on its face that it could be understood only as an effort to segregate voters into separate districts on the basis of race without sufficient justification. On remand, if the allegation of racial gerrymandering remained uncontradicted, the District Court was required to apply strict scrutiny---requiring the State to prove the plan was narrowly tailored to further a compelling governmental interest.

\section{Reasoning}
\textbf{1. Strict Scrutiny for Racial Classifications and Facial Bizarreness.} Justice O'Connor grounded the analysis in core Equal Protection principles: racial classifications are ``by their very nature odious to a free people'' and demand strict scrutiny regardless of whether they are labeled ``benign'' or ``remedial'' (\textit{Richmond v. J.A. Croson Co.}, 488 U.S. 469 (1989)). That rule extends to facially neutral laws that are ``unexplainable on grounds other than race'' (\textit{Arlington Heights v. Metropolitan Housing Dev. Corp.}, 429 U.S. 252 (1977)), including tortured voting boundaries (\textit{Gomillion v. Lightfoot}, 364 U.S. 339 (1960); \textit{Wright v. Rockefeller}, 376 U.S. 52 (1964)). Although legislatures are inevitably aware of race during redistricting, disregarding traditional districting criteria---compactness, contiguity, and respect for political subdivisions---provides objective evidence that a district was drawn solely to segregate voters by race.

\textbf{2. Expressive and Representational Harms (``Political Apartheid'').} Rejecting the dissenters' insistence that plaintiffs must show vote dilution, the majority identified distinct constitutional harms from racial sorting itself. Grouping geographically and politically separated citizens who share little but skin color ``bears an uncomfortable resemblance to political apartheid.'' It reinforces impermissible racial stereotypes that members of the same race think alike and prefer the same candidates, exacerbates racial bloc voting, and sends a ``pernicious'' message to elected officials that their primary duty is to represent one racial group rather than their entire constituency.

\textbf{3. Distinguishing \textit{UJO} and Political Gerrymanders.} The Court distinguished \textit{UJO} because the Hasidic plaintiffs there alleged vote dilution against a plan that adhered to traditional districting principles, not an irregular racial segregation claim. Likewise, racial gerrymanders are not governed by the relaxed standard applied to political gerrymanders (\textit{Davis v. Bandemer}), because the Fourteenth Amendment reserves its strictest scrutiny for racial classifications.

\textbf{4. Narrow Tailoring and VRA Compliance on Remand.} While acknowledging a State's strong interest in complying with federal antidiscrimination law, the Court cautioned that \S~5's nonretrogression rule does not give covered jurisdictions ``carte blanche'' to engage in racial gerrymandering beyond what is reasonably necessary. The Court left open on remand whether VRA \S~2 or remedying past discrimination could justify District 12 under strict scrutiny.

\section{Concurring and Dissenting Opinions}
\begin{itemize}
  \item \textbf{Justice White (Dissenting, joined by Justices Blackmun and Stevens):} Argued appellants alleged no cognizable constitutional injury because white voters remained a majority in 10 of 12 districts (83\% of seats vs. 79\% of voting-age population). Reading \textit{UJO}, \textit{Gomillion}, and \textit{Wright} as requiring both discriminatory purpose and dilutive effect, White contended that focusing on ``superficialities of shape and size'' invented an unworkable cause of action divorced from electoral harm, while obstructing legitimate VRA compliance.
  \item \textbf{Justice Blackmun (Dissenting):} Joined Justice White and underscored the irony of abandoning settled vote-dilution law to strike down the very plan under which North Carolina sent its first Black representatives to Congress since Reconstruction.
  \item \textbf{Justice Stevens (Dissenting):} Maintained that bizarre shapes merely evidence legislative purpose, which was already undisputed. Under the Equal Protection Clause, the duty to govern impartially is violated when a dominant majority distorts boundaries to entrench its own power, not when it facilitates representation for an underrepresented minority---whether defined by profession, religion, or race.
  \item \textbf{Justice Souter (Dissenting):} Emphasized two differences between redistricting and contracting or layoffs: race consciousness is unavoidable under the VRA, and assigning a voter to a district denies no one the franchise. Without dilutive harm, he saw no justification for untethering racial gerrymandering from vote-dilution doctrine solely because a district looks irregular.
\end{itemize}

\section{Key Quotations}
\begin{quote}
``Put differently, we believe that reapportionment is one area in which appearances do matter. A reapportionment plan that includes in one district individuals who belong to the same race, but who are otherwise widely separated by geographical and political boundaries, and who may have little in common with one another but the color of their skin, bears an uncomfortable resemblance to political apartheid.'' --- Justice O'Connor, 509 U.S. at 647 (\href{https://supreme.justia.com/cases/federal/us/509/630/}{Justia}).
\end{quote}
\begin{quote}
``When a district obviously is created solely to effectuate the perceived common interests of one racial group, elected officials are more likely to believe that their primary obligation is to represent only the members of that group, rather than their constituency as a whole. This is altogether antithetical to our system of representative democracy.'' --- Justice O'Connor, 509 U.S. at 648 (\href{https://supreme.justia.com/cases/federal/us/509/630/}{Justia}).
\end{quote}
\begin{quote}
``Racial gerrymandering, even for remedial purposes, may balkanize us into competing racial factions; it threatens to carry us further from the goal of a political system in which race no longer matters---a goal that the Fourteenth and Fifteenth Amendments embody, and to which the Nation continues to aspire.'' --- Justice O'Connor, 509 U.S. at 657 (\href{https://supreme.justia.com/cases/federal/us/509/630/}{Justia}).
\end{quote}
\begin{quote}
``Nonetheless, the notion that North Carolina's plan, under which whites remain a voting majority in a disproportionate number of congressional districts, and pursuant to which the State has sent its first black representatives since Reconstruction to the United States Congress, might have violated appellants' constitutional rights is both a fiction and a departure from settled equal protection principles.'' --- Justice White, dissenting, 509 U.S. at 660 (\href{https://supreme.justia.com/cases/federal/us/509/630/}{Justia}).
\end{quote}

\section{Doctrinal Significance and Modern Legacy (Through 2026)}
\textit{Shaw v. Reno} launched the modern ``racial gerrymandering'' Equal Protection doctrine, placing state legislatures in a recurring bind between the Voting Rights Act and the Fourteenth Amendment:
\begin{itemize}
  \item \textbf{From Bizarreness to Predominance:} Two years later, in \href{https://supreme.justia.com/cases/federal/us/515/900/}{\textit{Miller v. Johnson}, 515 U.S. 900 (1995)}, the Court clarified that bizarre shape is not a strict threshold; strict scrutiny applies whenever race is the ``predominant factor'' motivating district lines such that the legislature subordinates traditional race-neutral criteria to racial considerations. On remand in \href{https://supreme.justia.com/cases/federal/us/517/899/}{\textit{Shaw v. Hunt}, 517 U.S. 899 (1996)} (\textit{Shaw II}), the Court struck down North Carolina's District 12 under strict scrutiny.
  \item \textbf{Disentangling Race from Politics:} In \textit{Cooper v. Harris}, 581 U.S. 285 (2017), the Court struck down two North Carolina districts (including a later District 12) as racial gerrymanders, holding that plaintiffs must ``disentangle race from politics'' when race and party correlate, and that VRA \S~2 cannot justify race-predominant line-drawing without a ``strong basis in evidence.'' After \textit{Rucho v. Common Cause}, 588 U.S. 684 (2019), declared partisan gerrymandering nonjusticiable, \href{https://supreme.justia.com/cases/federal/us/602/22-807/}{\textit{Alexander v. South Carolina State Conference of the NAACP}, 602 U.S. 1 (2024)}, raised the evidentiary bar for \textit{Shaw} plaintiffs: courts must presume legislative good faith when a State asserts partisan aims, and a challenger's failure to submit an alternative map meeting the State's partisan goals with greater racial balance triggers an adverse inference.
  \item \textbf{\textit{Allen v. Milligan} (2023) to \textit{Louisiana v. Callais} (2026):} In \href{https://supreme.justia.com/cases/federal/us/599/21-1086/}{\textit{Allen v. Milligan}, 599 U.S. 1 (2023)}, a 5--4 Court reaffirmed the \textit{Thornburg v. Gingles} VRA \S~2 framework against Alabama's ``race-neutral benchmark'' challenge, though Justice Kavanaugh's concurrence questioned whether race-based redistricting under \S~2 could continue indefinitely. That tension reached a watershed in \href{https://supreme.justia.com/cases/federal/us/608/24-109/}{\textit{Louisiana v. Callais}, 608 U.S. \_\_\_ (2026)} (decided April 29, 2026, after reargument in October 2025). In a 6--3 opinion by Justice Alito (Thomas, J., concurring joined by Gorsuch; Kagan, J., dissenting joined by Sotomayor and Jackson), the Court struck down Louisiana's SB8---whose second majority-Black district (District 6) snaked roughly 250 miles from Shreveport to Baton Rouge---as an unconstitutional racial gerrymander. Resolving the question left open since \textit{Shaw} and \textit{Miller}, \textit{Callais} held that VRA \S~2 compliance supplies a compelling interest only when \S~2 is construed to enforce the Fifteenth Amendment's ban on intentional discrimination (requiring a strong inference of racial purpose), updated \textit{Gingles} to bar race-based illustrative maps and require plaintiffs to control for partisan affiliation and meet all legitimate state political/incumbency goals, and concluded SB8 failed strict scrutiny.
\end{itemize}

\section{Pedagogical Critiques and Exam Hypotheticals}
\textbf{Pedagogical Critiques:}
\begin{itemize}
  \item \textbf{Standing and Expressive Harm:} Unlike traditional Equal Protection plaintiffs denied a job or seat, \textit{Shaw} plaintiffs suffer no vote dilution. Critics argue the Court's ``expressive harm'' theory turns aesthetic distaste into a constitutional injury, whereas defenders contend that state-sponsored racial sorting itself violates individual dignity.
  \item \textbf{The Race--Partisanship Paradox:} Together, \textit{Shaw}, \textit{Rucho}, \textit{Alexander}, and \textit{Callais} create a stark asymmetry: sorting Black voters \textit{because they are Black} triggers near-fatal strict scrutiny, yet sorting the exact same voters \textit{because they vote Democratic} is constitutionally permissible and federally nonjusticiable.
\end{itemize}

\textbf{Exam Hypothetical:}
After the 2030 census, a State with a 28\% Hispanic voting-age population enacts a congressional map containing a 210-mile, barbell-shaped District 4 that splits nine counties to connect two distant urban Hispanic enclaves, achieving a 52\% Hispanic voting-age population. Legislative emails show mapmakers set a ``50\%+ HVAP floor'' to preempt a threatened VRA \S~2 lawsuit, though no court had found intentional discrimination and no illustrative map could create the district while preserving the State's partisan-advantage goals. A white voter residing in District 4 sues under the Equal Protection Clause.
\begin{itemize}
  \item \textbf{\textit{Shaw} / \textit{Miller} / \textit{Callais} Analysis:} The plaintiff has standing as a District 4 resident (\textit{Shaw II}) and states a strong \textit{Shaw} claim. Both the bizarre, county-splitting geometry and the direct email evidence of a 50\%+ racial target show that race predominated over traditional districting principles, triggering strict scrutiny. Under \textit{Louisiana v. Callais} (2026), the State cannot establish a compelling interest in VRA \S~2 compliance because \S~2 requires a strong inference of intentional racial discrimination and an illustrative map that satisfies the State's legitimate political goals without using race; District 4 is therefore an unconstitutional racial gerrymander.
  \item \textbf{``Opposite Side'' Contrast:} Now alter the facts: suppose the legislature draws a compact, county-respecting District 4 that moves 20,000 heavily Democratic Hispanic voters \textit{out} of District 4 into neighboring District 5 to lock in a safe Republican incumbent in District 4, without any racial target in the legislative record, and challengers fail to produce an alternative map that achieves the same Republican partisan margin while keeping those voters in District 4. Under \textit{Alexander v. South Carolina State Conference of the NAACP} (2024) and \textit{Rucho}, the presumption of legislative good faith and the adverse inference from the missing alternative map defeat the \textit{Shaw} racial-gerrymandering claim because the challengers failed to disentangle race from permissible partisan politics.
\end{itemize}

\end{document}
